% CHAPTER 2

\chapter{Theory}
\label{chap:theory}

% 2.1 EXACT HYDROGENIC GROUND STATE

The numerical investigation undertaken requires a theoretical framework against
which the Gaussian approximation can be assessed. This chapter begins
with the exact hydrogenic ground state, which provides a benchmark
for evaluating the approximate solutions. The variational formulation
and analytical Gaussian matrix elements are then developed, followed
by the measures of accuracy, numerical conditioning and nuclear-charge
scaling used throughout the subsequent analysis.

\section{Exact Hydrogenic Ground State}

In atomic units
($\hbar=m_e=e=4\pi\varepsilon_0=1$),
the Hamiltonian for an electron bound to a fixed
point nucleus of charge $+Ze$ is
\begin{equation}
    \hat H=-\frac12\nabla^2-\frac{Z}{r},
    \label{eq:hydrogen_hamiltonian}
\end{equation}
where the two terms represent the electron's kinetic
energy and its attractive Coulomb interaction with
the nucleus. The exact normalised ground-state
wavefunction and corresponding energy are
\cite{bransden_physics_2009}
\begin{equation}
    \psi_{1s}(r)
    =\left(\frac{Z^3}{\pi}\right)^{1/2}e^{-Zr},
    \qquad
    E_{1s}=-\frac{Z^2}{2}.
    \label{eq:exact_state}
\end{equation}

The probability density $|\psi(r)|^2$ gives the
probability per unit volume of finding the electron
at a particular position. For a spherically
symmetric state, the radial probability density is
\begin{equation}
    P(r)=4\pi r^2|\psi(r)|^2,
    \qquad
    \int_0^\infty P(r)\,\mathrm dr=1.
    \label{eq:radial_probability}
\end{equation}
The factor $4\pi r^2$ accounts for the volume of a
spherical shell. Substituting the exact $1s$
wavefunction gives
\[
    P_{1s}(r)=4Z^3r^2e^{-2Zr}.
\]
Although the wavefunction is largest at the nucleus,
the radial probability reaches its maximum at
$r=1/Z$. The mean radius is
$\langle r\rangle=3/(2Z)$, reflecting the contribution
from the full radial distribution rather than only
its most probable radius.

These exact results provide complementary benchmarks
for the Gaussian approximation. The ground-state
energy measures the quality of the variational result,
while the wavefunction, radial probability and mean
radius describe its spatial accuracy. Both are needed
because an accurate energy does not necessarily
imply an equally accurate wavefunction.



% 2.2 GAUSSIAN BASIS AND VARIATIONAL METHOD
\newpage
\section{Gaussian Basis and Variational Method}

The approximate ground-state wavefunction is
represented by a linear combination of $N$
spherical Gaussian functions centred on the
nucleus \cite{quiney_art_2018},
\begin{equation}
    \psi_{\mathrm G}(r)
    =\sum_{i=1}^{N}c_i\phi_i(r),
    \qquad
    \phi_i(r)=A_i e^{-\lambda_i r^2},
    \qquad
    A_i=\left(\frac{2\lambda_i}{\pi}\right)^{3/4}.
    \label{eq:gaussian_expansion}
\end{equation}
The coefficients $c_i$ determine the contribution
of each function, while $A_i$ normalises each
Gaussian individually. The exponents follow an
even-tempered sequence,
\begin{equation}
    \lambda_i=\alpha\beta^{i-1},
    \qquad
    \alpha>0,\quad\beta>1,\quad i=1,\ldots,N.
    \label{eq:exponent_sequence}
\end{equation}
Here $\alpha$ sets the smallest exponent and
$\beta$ determines the ratio between successive
exponents
\cite{quiney_art_2018,BaR,FellerRuedenberg1979,Wilson}.

The characteristic width of each Gaussian is
proportional to $1/\sqrt{\lambda_i}$. Small
exponents describe more diffuse regions of the
electron distribution, while large exponents
produce narrower functions concentrated near
the nucleus. Since the largest exponent is
$\alpha\beta^{N-1}$, increasing $N$ at fixed
$\alpha$ and $\beta$ adds progressively narrower
functions without extending the diffuse end of
the basis. Consequently, increasing the number
of functions does not necessarily improve the
description of every radial region equally.

Gaussian functions are particularly useful
because their products and many of the required
integrals can be evaluated analytically.
The Gaussian product theorem expresses the
product of two Gaussians as another Gaussian,
simplifying the construction of the matrices
used in electronic-structure calculations
\cite{quiney_art_2018}. Although the exact
hydrogenic wavefunction decays as $e^{-Zr}$
rather than $e^{-\lambda r^2}$, a suitable
combination of Gaussians can approximate its
spatial behaviour while retaining these
computational advantages.

The expansion coefficients are obtained by
minimising the expectation value of the
Hamiltonian. For a nonorthogonal basis,
the energy is expressed through the
Rayleigh quotient,
\begin{equation}
    E[\mathbf c]
    =\frac{\mathbf c^{\mathsf T}\mathbf H\mathbf c}
           {\mathbf c^{\mathsf T}\mathbf S\mathbf c},
    \qquad
    H_{ij}=\langle\phi_i|\hat H|\phi_j\rangle,
    \quad
    S_{ij}=\langle\phi_i|\phi_j\rangle.
    \label{eq:rayleigh_quotient}
\end{equation}
The Hamiltonian matrix $\mathbf H$ contains
the kinetic and potential energy contributions,
while the overlap matrix $\mathbf S$ accounts
for the nonorthogonality of the Gaussian
basis functions.

According to the variational principle, the
expectation value of the Hamiltonian for any
normalised trial wavefunction cannot be lower
than the exact ground-state energy
\cite{bransden_physics_2009}. The coefficients
are therefore chosen to minimise the energy
subject to
$\mathbf c^{\mathsf T}\mathbf S\mathbf c=1$.
Introducing the Lagrange multiplier $E$,
identified with the stationary energy, gives
\[
    \mathcal L
    =\mathbf c^{\mathsf T}\mathbf H\mathbf c
    -E\left(
       \mathbf c^{\mathsf T}\mathbf S\mathbf c-1
      \right).
\]
Differentiating with respect to the coefficient
vector,
\[
    \frac{\partial\mathcal L}{\partial\mathbf c}
    =2\mathbf H\mathbf c
     -2E\mathbf S\mathbf c=0,
\]
leads to the generalised eigenvalue equation
\begin{equation}
    \mathbf H\mathbf c=E\mathbf S\mathbf c.
    \label{eq:generalised_eigenproblem}
\end{equation}
This formulation is also used in the
computational quantum-chemistry approach
described by Quiney and Flynn
\cite{quiney_art_2018}. Further details
are given in Appendix~\ref{app:gaussian-integrals}.

The overlap matrix appears because the
Gaussian functions are not mutually
orthogonal, so $\mathbf S\neq\mathbf I$.
For a linearly independent basis,
$\mathbf S$ is positive definite and
the generalised eigenvalues are real.
The lowest eigenvalue gives the minimum
energy available within the chosen
Gaussian space. The variational bound
therefore requires
$E_{\mathrm G}\geq E_{1s}$ in exact
arithmetic, providing a useful check
on the numerical solution.

The matrix elements can be evaluated
analytically. Defining
$p_{ij}=\lambda_i+\lambda_j$, the overlap,
kinetic-energy and Coulomb potential
matrices are
\begin{align}
    S_{ij}
    &=A_iA_j
      \left(\frac{\pi}{p_{ij}}\right)^{3/2},
    \label{eq:overlap_matrix}\\
    T_{ij}
    &=\frac{3\lambda_i\lambda_j}{p_{ij}}S_{ij},
    \label{eq:kinetic_matrix}\\
    V_{ij}
    &=-\frac{2\pi ZA_iA_j}{p_{ij}},
    \label{eq:potential_matrix}\\
    H_{ij}
    &=T_{ij}+V_{ij}.
    \label{eq:hamiltonian_matrix}
\end{align}
These expressions agree with the
one-centre Gaussian matrix elements
described by Quiney and Flynn
\cite{quiney_art_2018}. Their derivations,
including Gaussian normalisation, are
provided in Appendix~\ref{app:gaussian-integrals}.

Evaluating the matrix elements analytically
avoids numerical integration errors in
the construction of $\mathbf H$ and
$\mathbf S$. It also makes repeated
calculations with different basis sizes
and exponent parameters more efficient.
The resulting generalised eigenvalue
problem is solved numerically using the
methods described in Chapter~3.


% 2.3 ACCURACY, CONDITIONING AND SCALING

\section{Accuracy, Conditioning and Nuclear-Charge Scaling}

The primary measure of accuracy is the
relative ground-state energy error,
\begin{equation}
    \varepsilon_E
    =\frac{|E_{\mathrm G}-E_{1s}|}{|E_{1s}|}.
    \label{eq:energy_error}
\end{equation}
However, the variational energy alone
does not establish how closely the
approximate wavefunction reproduces
the exact state. Spatial accuracy is
therefore also assessed through the
relative error in the mean radius,
\begin{equation}
    \varepsilon_r
    =\frac{|\langle r\rangle_{\mathrm G}-3/(2Z)|}
           {3/(2Z)},
    \qquad
    \langle r\rangle_{\mathrm G}
    =\int_0^\infty rP_{\mathrm G}(r)\,\mathrm dr,
    \label{eq:radius_error}
\end{equation}
and the three-dimensional wavefunction
error, expressed as an $L^2$ norm,
\begin{equation}
    \varepsilon_\psi
    =\left[
        4\pi\int_0^\infty r^2
        |\psi_{\mathrm G}(r)-\psi_{1s}(r)|^2
        \,\mathrm dr
      \right]^{1/2}.
    \label{eq:wavefunction_error}
\end{equation}
Before calculating this difference,
the overall sign of the numerical
wavefunction is chosen to agree
with the positive exact ground state.

The variational principle explains
why the energy may be substantially
more accurate than the wavefunction.
For a normalised Gaussian approximation,
write
\[
    \psi_{\mathrm G}
    =a\psi_{1s}+\delta\psi,
    \qquad
    \langle\psi_{1s}|\delta\psi\rangle=0.
\]
The exact ground state satisfies
$\hat H\psi_{1s}=E_{1s}\psi_{1s}$.
Using this relation and the
normalisation of both states,
the energy difference becomes
\begin{equation}
    E_{\mathrm G}-E_{1s}
    =\langle\delta\psi|
      (\hat H-E_{1s})|\delta\psi\rangle.
    \label{eq:variational_error}
\end{equation}
The terms linear in the wavefunction
difference cancel, leaving an
energy-weighted quadratic expression.
A small energy error therefore does
not necessarily imply a comparably
small wavefunction error. This
distinction motivates the separate
accuracy measures examined in
Chapter~4.

Another limitation comes from the
behaviour of the wavefunction at
the nucleus. The exact hydrogenic
wavefunction satisfies the
electron--nucleus cusp condition,
whereas a finite sum of centred
spherical Gaussians has zero radial
derivative at the origin:
\begin{equation}
    \psi_{1s}'(0)=-Z\psi_{1s}(0),
    \qquad
    \psi_{\mathrm G}'(0)=0.
    \label{eq:cusp_condition}
\end{equation}
The cusp follows from the singular
Coulomb potential. For a spherically
symmetric wavefunction,
\[
    \nabla^2\psi
    =\psi''(r)+\frac{2}{r}\psi'(r).
\]
If $\psi'(0)\neq0$, the second term
contains a contribution proportional
to $1/r$. For the exact hydrogenic
state, the cusp condition ensures
that this singular kinetic-energy
contribution cancels the Coulomb
potential as $r\to0$. A finite
Gaussian expansion cannot reproduce
this cancellation exactly
\cite{Kato,Kutzelnigg1994}.

The effect can be examined through
the local energy,
\begin{equation}
    E_{\mathrm{loc}}(r)
    =\frac{\hat H\psi(r)}{\psi(r)}.
    \label{eq:local_energy}
\end{equation}
For an exact eigenfunction, the
local energy equals the corresponding
eigenvalue wherever the wavefunction
is nonzero. For a finite Gaussian
expansion with nonzero amplitude
at the nucleus, the kinetic-energy
contribution remains finite as
$r\to0$, while the Coulomb term
diverges as $-Z/r$. The local
energy consequently tends towards
negative infinity near the nucleus,
revealing a limitation that may
be less apparent from the total
variational energy.

Increasing the number of Gaussians
can improve the approximation,
but it can also introduce numerical
difficulties when basis functions
become nearly linearly dependent.
This is assessed using the spectral
condition number of the overlap
matrix,
\begin{equation}
    \kappa(\mathbf S)
    =\frac{s_{\max}}{s_{\min}},
    \label{eq:condition_number}
\end{equation}
where $s_{\max}$ and $s_{\min}$
are its largest and smallest
eigenvalues. For a positive-definite
overlap matrix, both are positive.
A large condition number indicates
that some combinations of basis
functions are nearly redundant,
making the generalised eigenvalue
problem more sensitive to numerical
rounding errors \cite{Lehtola2019}.

It is therefore important to
distinguish between an incomplete
basis and an ill-conditioned one.
An incomplete basis may fail to
represent the physical wavefunction
accurately even when the matrix
problem is solved correctly.
Conversely, a larger basis may
contain nearly dependent functions
that make the numerical solution
unreliable. In the parameter
search, $\kappa(\mathbf S)\leq10^{12}$
is used as a practical screening
criterion rather than a universal
guarantee of accuracy.

Hydrogenic wavefunctions also have
a simple dependence on nuclear
charge. The exact ground state
depends on the scaled coordinate
$\rho=Zr$, so its characteristic
radial length decreases as $1/Z$
\cite{bransden_physics_2009}.
To preserve the Gaussian widths
relative to this length scale,
the exponents must satisfy
\begin{equation}
    \lambda_i(Z)=Z^2\lambda_i(1).
    \label{eq:scaled_exponents}
\end{equation}
For an even-tempered basis,
this corresponds to replacing
$\alpha$ with $Z^2\alpha$ while
keeping $N$ and $\beta$ unchanged.

The exponent scaling follows
by writing each Gaussian in
terms of $\rho$:
\[
    e^{-\lambda_i(Z)r^2}
    =e^{-[\lambda_i(Z)/Z^2]\rho^2}.
\]
Choosing $\lambda_i(Z)=Z^2\lambda_i(1)$
therefore preserves the Gaussian
shape in the scaled coordinate.

The corresponding energy scaling
can be obtained directly from the
Hamiltonian. Since
$\nabla_r^2=Z^2\nabla_\rho^2$
and $r=\rho/Z$,
\[
    \hat H(Z)
    =Z^2\left(
        -\frac12\nabla_\rho^2-\frac{1}{\rho}
      \right).
\]
Thus, for a consistently scaled
Gaussian basis, the variational
energy satisfies
\begin{equation}
    E_{\mathrm G}(Z)=Z^2E_{\mathrm G}(1).
    \label{eq:scaled_energy}
\end{equation}
Because the exact energy also
scales as $Z^2$, the relative
energy error is unchanged under
this transformation.

The normalised wavefunction scales
as $\psi_Z(r)=Z^{3/2}\psi_1(Zr)$,
giving the corresponding radial
probability relation,
\begin{equation}
    P_Z(r)=ZP_1(Zr),
    \qquad
    \frac{P_Z(\rho/Z)}{Z}=P_1(\rho).
    \label{eq:radial_scaling}
\end{equation}
The scaled radial distributions
should therefore coincide when
expressed in terms of $\rho=Zr$.
Together, these relations provide
specific theoretical predictions
for the nuclear-charge calculations
presented in Chapter~4.
